% ============================================================
%  Section 6 — Discussion; Section 7 — Conclusion
% ============================================================

\section{Discussion}
\label{sec:discussion}

\subsection{Role of Each Contribution}

The three contributions do not act in the same way. Multi-task learning
(C2) is the only component that improves the point prediction, and its
effect is largest on the healthy windows, where the stage task gives
the encoder an explicit signal that the RUL target is flat. Its role
does not end with accuracy: its stage head supplies the grouping that
makes the conformal intervals of C3 about 20\% narrower than a single
quantile over the whole life, at nominal coverage. The auxiliary task
therefore serves twice, first to shape the representation and then to
condition the uncertainty, and a single network provides the RUL, the
stage and the interval.\\

The calibration scheme of C3 mattered more than any modelling choice
for the validity of the intervals. With two calibration bearings, no
version of the intervals reached its target, including the version
grouped by the true stage. With CV+, every version came close to it.
On a dataset of fifteen units, the number of bearings that inform the
calibration is the binding constraint, and a held-out calibration split
spends too few of them. The price is computational: CV+ trains $K$
models per fold instead of one.\\

Generative augmentation (C1) produced windows that are hard to
distinguish from real ones and as diverse, but it did not change
accuracy or interval quality beyond the seed-to-seed variability, alone
or combined with C2. The duplication control points to the reason: the
predictor did not improve when the number of real near-failure windows
was increased either, so it was not limited by their number. Earlier
GAN studies on XJTU-SY report gains of 15--30\% in
RMSE~\cite{9486650, lu_joint_2022}, but on scalar health indicators and
with other targets and protocols; with a multivariate input, a
piecewise target and multi-task training, the benefit of additional
near-failure windows was not measurable here. The finding is negative
but informative, and it is reported as such rather than as a gain.

\subsection{What the Guarantee Covers}

The coverage guarantee of conformal prediction is marginal: it holds on
average over bearings drawn from the same population as the
calibration bearings, not for every bearing. The per-bearing analysis
makes this concrete. Most bearings are covered at or above the target,
but a few are covered far below it, and three of the four weakest are
among the shortest-lived. On those bearings the predictor is wrong
rather than merely uncertain: Bearing~1-4 receives narrow intervals
that miss most of its near-failure windows. No calibration scheme can
repair such errors, as the bearing-balanced variant confirms, because
they come from degradation behaviour that the other bearings do not
exhibit. Reporting coverage per bearing is therefore necessary to
interpret an aggregate coverage on a dataset of this size.\\

Conditioning on the predicted stage also moves coverage between stages.
The guarantee holds within the predicted stages, and on the true stages
the near-failure windows are over-covered while the early-degradation
windows are under-covered. The intervals near failure are about 35\%
wider than they would be with correct stage recognition. Improving the
separation between early degradation and near failure is therefore the
most direct way to narrow the intervals where decisions are taken. From
a maintenance standpoint the current direction of error is the safer
one: near failure, the intervals are conservative rather than
optimistic.

\subsection{Implications for Maintenance Decisions}

The interval supports decisions that a point estimate does not. A
maintenance rule can act on the lower bound of the interval, scheduling
intervention when the lower bound of the RUL falls below the lead time
needed to plan it, so that the decision is protected against
overestimation of the remaining life at the chosen confidence level.
Because the interval is conditioned on the predicted stage, the same
rule automatically becomes more cautious once a bearing is recognised
as near failure. The per-bearing results also indicate when such a rule
should not be trusted alone: on short-lived bearings that fail
abruptly, a confident prediction can still be wrong, and a fixed
inspection schedule or a threshold on the raw health indicator remains
a necessary safeguard.

\subsection{Limitations and Future Work}

The study uses a single dataset of fifteen bearings, and its
conclusions should be confirmed on other run-to-failure datasets, such
as PRONOSTIA and IMS, and on other components. The stage labels are
derived from a health-indicator criterion that uses the whole
trajectory; they are consistent with the visible degradation but are
not ground truth, and the reference results by true stage inherit this
definition. CV+ multiplies the training cost by the number of inner
groups, including one generator per inner model. Finally, the intervals
are computed independently for each window; for a bearing in service,
an adaptive conformal scheme that updates the calibration as new
recordings arrive, and a grouping based on a better recognition of the
transition to near failure, are natural extensions.

% ------------------------------------------------------------
\section{Conclusion}
\label{sec:conclusion}

This paper proposed a three-stage framework for bearing RUL prediction
combining stage-conditioned generative augmentation, multi-task RUL and
stage prediction, and stage-conditional cross-conformal intervals, with
stage labels and RUL targets derived automatically from a guarded
health-indicator criterion. On XJTU-SY, with bearing-level
cross-validation and five seeds, multi-task learning lowered the
near-failure RMSE from 0.225 to 0.206 and the overall RMSE from 0.385
to 0.291. Cross-conformal calibration on all non-test bearings restored
nominal coverage, 0.917 overall at a 90\% target, where split
calibration on two bearings reached at most 0.83. Conditioning the
intervals on the stage predicted by the multi-task model narrowed them
by about 20\% at higher coverage, so that a single network delivers the
RUL estimate, the degradation stage and a calibrated interval.
Generative augmentation produced realistic near-failure windows but did
not improve accuracy or intervals measurably. Coverage held on average
over bearings but not on every bearing, with the shortest-lived
bearings the least well covered, which shows that per-bearing coverage
should accompany aggregate coverage when conformal intervals are
reported on small run-to-failure datasets.
